\uuid{ZIua}
\chapitre{Probabilité discrète}
\niveau{L2}
\module{Probabilité}
\sousChapitre{Variable aléatoire discrète}
\cpge{mpsi-mp2i.probabilites-univers-fini.d, pcsi.probabilites-variables-aleatoires.d, ptsi.probabilites-variables-aleatoires.d, mpsi-mp2i.esperance-variance.b, pcsi.esperance-variance.b, ptsi.esperance-variance.b, bcpst1.variable-aleatoire-univers-fini.a, bcpst1.variable-aleatoire-univers-fini.c}
\titre{Loi d'une variable aléatoire}
\theme{variables aléatoires discrètes}
\auteur{}
\datecreate{2023-02-07}
\organisation{AMSCC}
\difficulte{}
\contenu{

\texte{ 	Une variable aléatoire $X$ prend les valeurs 3,4,5 et 6. }
	\begin{enumerate}
		\item \question{ Quelle est la loi de probabilité de X sachant que
		$$\PP(X > 5) = 1/2 \qquad \PP(X < 5) = 1/3 \qquad P(X=3)=P(X=4)$$ }
		\item \question{ Calculer l'espérance mathématique, la variance puis l'écart type de la variable $X$. }
	\end{enumerate}}
